\subsection{Equicontinuity criteria of multilinear mappings for topologies defined by semi-norms}

PROPOSITION 4. --- Let \(E_{i}\) ( \(1 \leqslant i \leqslant n\) ) and F be topological vector spaces over K; we suppose that, for every i, the topology of \(E_{i}\) is defined by a directed set of semi-norms \(\Gamma_{i}\) , and that the topology of F is defined by a set of semi-norms \(\Gamma\) . Then a set H, of multilinear mappings of \(\prod_{i=1}^{n} \mathrm{E}_{i}\) in F is equicontinuous if, and only if, for each semi-norm \(q \in \Gamma\) , and each index \(i\) , there exists a semi-norm \(p_{i} \in \Gamma_{i}\) , and a number \(a > 0\) , such that for each function \(u \in \mathrm{H}\) and point \((x_{i}) \in \prod_{i=1}^{n} \mathrm{E}_{i}\) ,

\[
q \left(u \left(x _ {1}, x _ {2}, \dots , x _ {n}\right)\right) \leqslant a. p _ {1} \left(x _ {1}\right) p _ {2} \left(x _ {2}\right) \dots p _ {n} \left(x _ {n}\right).\tag{5}
\]

The condition is sufficient since it implies that H is equicontinuous at \((0, 0, \ldots, 0)\) and therefore everywhere (I, p.~9, prop. 6).

We show that the condition is necessary. By hypothesis, for every semi-norm \(q \in \Gamma\) and every number \(\beta > 0\) , we have \(q(u(x_1, x_2, ..., x_n)) \leqslant \beta\) for every function \(u \in \mathbf{H}\) provided that \(p_i(x_i) \leqslant \alpha_i\) are true for each index \(i\) , \(1 \leqslant i \leqslant n\) , and certain appropriately chosen numbers \(\alpha_i > 0\) and semi-norms \(p_i \in \Gamma_i\) . As \(\mathbf{K}\) is non-discrete, we can also suppose that, for every \(i\) , we have \(\alpha_i = |\lambda_i| < 1\) where \(\lambda_i \in \mathbf{K}\) . Then let \((x_1, x_2, ..., x_n)\) be any point of \(\prod_{i=1}^{n} E_i\) , and for each index \(i\) , let \(m_i \in \mathbf{Z}\) be an integer such that \(p_i(x_i) \leqslant |\lambda_i|^{m_i + 1}\) ; this can be written as \(p_i(\lambda_i^{-m_i} x_i) \leqslant |\lambda_i| (1 \leqslant i \leqslant n)\) , therefore, by hypothesis, we have

\[
q \left(u \left(x _ {1}, x _ {2}, \dots , x _ {n}\right)\right) \leqslant \beta \left| \lambda_ {1} \right| ^ {m _ {1}} \left| \lambda_ {2} \right| ^ {m _ {2}} \dots \left| \lambda_ {n} \right| ^ {m _ {n}}.\tag{6}
\]

Suppose firstly that one of the \(p_{i}(x_{i})\) is zero, then we can take \(m_{i} \in N\) arbitrarily large, therefore \(q(u(x_{1}, x_{2}, ..., x_{n})) = 0\) . If, on the contrary, all the \(p_{i}(x_{i})\) are \(\neq 0\) , take the integer \(m_{i}\) such that \(|\lambda_{i}|^{m_{i}+2} < p_{i}(x_{i}) \leqslant |\lambda_{i}|^{m_{i}+1}\) for each i; then we have \(|\lambda_{i}|^{m_{i}} < |\lambda_{i}|^{-2}p_{i}(x_{i})\) , from which, by (6), the relation (5) follows with

\[
a = \beta (| \lambda_ {1} | \cdot | \lambda_ {2} | \dots | \lambda_ {n} |) ^ {- 2}. \qquad \text{Q.E.D.}
\]

COROLLARY. --- The set H is equicontinuous if, and only if, for every semi-norm \(q \in \Gamma\) , there exists a neighbourhood of 0 in \(\prod_{i=1}^{n} E_{i}\) , in which the functions \(q \circ u\) , for \(u \in H\) , are uniformly bounded.

The condition is evidently necessary, and the demonstration of prop. 4 shows that it implies an inequality of the form (5) for all \(u \in H\) , and therefore the equicontinuity of H.

We state explicitly the particular case of prop. 4 for linear mappings.

PROPOSITION 5. --- Let E, F be two topological vector spaces over a non-discrete valued division ring K; suppose that the topology of E (resp. F) is defined by a set of semi-norms \(\Gamma\) (resp. \(\Gamma'\) ). Let H be a set of linear mappings of E in F. The following conditions are equivalent :

\begin{enumerate}
\def\labelenumi{\alph{enumi})}
\item
  H is equicontinuous.
\item
  For every semi-norm \(q \in \Gamma'\) , there exists a finite family \((p_i)_{1 \leq i \leq n}\) of semi-norms belonging to \(\Gamma\) and a number \(a > 0\) such that, for all \(x \in E\) and all \(u \in H\) ,
\end{enumerate}

\[
q (u (x)) \leqslant a. \sup _ {1 \leqslant i \leqslant n} p _ {i} (x).\tag{7}
\]

\begin{enumerate}
\def\labelenumi{\alph{enumi})}
\setcounter{enumi}{2}
\tightlist
\item
  For every semi-norm \(q \in \Gamma'\) , the mapping \(\sup_{u \in H} (q \circ u)\) is a continuous semi-norm on \(E\) .
\end{enumerate}

COROLLARY 1. --- Suppose that \(\mathcal{T}, \mathcal{T}'\) are two topologies on a vector space \(\mathbf{E}\) over \(\mathbf{K}\) defined, respectively, by two sets of semi-norms \(\Gamma\) and \(\Gamma'\) . \(\mathcal{T}\) is finer than \(\mathcal{T}'\) if, and only if, for every semi-norm \(q \in \Gamma'\) , there exists a finite family \((p_i)_{1 \leqslant i \leqslant n}\) of semi-norms belonging to \(\Gamma\) and a number \(a > 0\) such that, for all \(x \in \mathbf{E}\) , we have \(q(x) \leqslant a . \sup_{1 \leqslant i \leqslant n} p_i(x)\) .

In fact this shows that the identity mapping of E with topology T, on E with topology \(T'\) , is continuous.

COROLLARY 2. --- Suppose that the topology \(\mathcal{T}\) of a topological vector space \(E\) over \(K\) is defined by a directed set of semi-norms \(\Gamma\) ; for each semi-norm \(p \in \Gamma\) , let \(E_p\) be the space obtained from \(E\) using the topology defined by \(p\) . The set \(E'\) of linear forms on \(E\) that are continuous for \(\mathcal{T}\) is the union of the sets \(E_p'\) , where \(E_p'\) is the set of continuous linear forms in \(E_p\) ( \(p \in \Gamma\) ).

